\documentclass[../../main.tex]{subfiles}
\begin{document}
\section{Lie 群同态}\label{section:ch07-Lie群同态}

若 $G$ 和 $H$ 是 Lie 群，则从 $G$ 到 $H$ 的一个\textbf{Lie 群同态}是一个同时也是群同态的光滑映射 $F\colon G\to H$。若它同时也是微分同胚，则称为\textbf{Lie 群同构}，这意味着它有一个逆映射，且该逆映射也是 Lie 群同态。此时我们称 $G$ 与 $H$ 是\textbf{同构的}Lie 群。

\begin{example}\label{example:ch07-04}[Lie 群同态]
\begin{enumerate}[(a)]
\item 包含映射 $S^1\hookrightarrow\mathbb{C}^*$ 是 Lie 群同态。

\item 将 $\mathbb{R}$ 视为加法下的 Lie 群，$\mathbb{R}^*$ 视为乘法下的 Lie 群，则映射 $\exp\colon\mathbb{R}\to\mathbb{R}^*$ 由 $\exp(t)=e^t$ 给出，它是光滑的，并且是 Lie 群同态，因为 $e^{s+t}=e^s e^t$。$\exp$ 的像为由正实数组成的开子群 $\mathbb{R}^+$，且 $\exp\colon\mathbb{R}\to\mathbb{R}^+$ 是一个 Lie 群同构，其逆为 $\log\colon\mathbb{R}^+\to\mathbb{R}$。

\item 类似地，$\exp\colon\mathbb{C}\to\mathbb{C}^*$ 由 $\exp(z)=e^z$ 给出，是一个 Lie 群同态。它是满射但不是单射，因为其核由形如 $2\pi\mathrm{i}k$ 的复数组成，其中 $k$ 为整数。

\item 映射 $\varepsilon\colon\mathbb{R}\to S^1$ 由 $\varepsilon(t)=e^{2\pi\mathrm{i}t}$ 定义，是一个 Lie 群同态，其核为整数集 $\mathbb{Z}$。类似地，映射 $\varepsilon^n\colon\mathbb{R}^n\to\mathbb{T}^n$ 由
\begin{align*}
\varepsilon^n(x^1,\dots,x^n)=\bigl(e^{2\pi\mathrm{i}x^1},\dots,e^{2\pi\mathrm{i}x^n}\bigr)
\end{align*}
定义，是一个 Lie 群同态，其核为 $\mathbb{Z}^n$。

\item 行列式函数 $\det\colon\operatorname{GL}(n,\mathbb{R})\to\mathbb{R}^*$ 是光滑的，因为 $\det A$ 是 $A$ 的矩阵元的多项式。它是 Lie 群同态，因为 $\det(AB)=(\det A)(\det B)$。类似地，$\det\colon\operatorname{GL}(n,\mathbb{C})\to\mathbb{C}^*$ 是 Lie 群同态。

\item 若 $G$ 是 Lie 群且 $g\in G$，则由 $g$ 的\textbf{共轭}是映射 $C_g\colon G\to G$，由 $C_g(h)=ghg^{-1}$ 给出。由于群乘法与取逆都是光滑的，$C_g$ 是光滑的，简单计算可知它是群同态。事实上，它是同构，因为 $C_{g^{-1}}$ 是其逆。子群 $H\subseteq G$ 称为\textbf{正规的}，若对每个 $g\in G$ 都有 $C_g(H)=H$。
\end{enumerate}
\end{example}

下一个定理对于理解 Lie 群同态的许多性质十分重要。

\begin{theorem}\label{theorem:ch07-05}
每个 Lie 群同态都具有常秩。
\end{theorem}

\begin{proof}
设 $F\colon G\to H$ 是 Lie 群同态，分别记 $e$ 和 $\tilde{e}$ 为 $G$ 和 $H$ 的单位元。设 $g_0$ 是 $G$ 的任意元素。我们将证明 $\mathrm{d}F_{g_0}$ 与 $\mathrm{d}F_e$ 具有相同的秩。$F$ 是同态意味着对所有 $g\in G$，
\begin{align*}
F\bigl(L_{g_0}(g)\bigr)=F(g_0g)=F(g_0)F(g)=L_{F(g_0)}\bigl(F(g)\bigr),
\end{align*}
换言之，$F\circ L_{g_0}=L_{F(g_0)}\circ F$。在单位元处对两边取微分并利用\cref{proposition:ch03-06}(b)，得到
\begin{align*}
\mathrm{d}F_{g_0}\circ\mathrm{d}(L_{g_0})_e=\mathrm{d}(L_{F(g_0)})_{\tilde{e}}\circ\mathrm{d}F_e.
\end{align*}
Lie 群中任意元素的左平移都是微分同胚，因此 $\mathrm{d}(L_{g_0})_e$ 和 $\mathrm{d}(L_{F(g_0)})_{\tilde{e}}$ 都是同构。由于与同构复合不改变线性映射的秩，由此可知 $\mathrm{d}F_{g_0}$ 与 $\mathrm{d}F_e$ 具有相同的秩。
\end{proof}

\begin{corollary}\label{corollary:ch07-06}
Lie 群同态是 Lie 群同构当且仅当它是双射。
\end{corollary}

\begin{proof}
整体秩定理表明，双射的 Lie 群同态是微分同胚。
\end{proof}

\textbf{万有覆叠群}

覆叠空间理论给出如下关于 Lie 群的重要结果。

\begin{theorem}\label{theorem:ch07-07}[万有覆叠群的存在性]
设 $G$ 是连通 Lie 群。存在一个单连通 Lie 群 $\widetilde{G}$，称为 $G$ 的\textbf{万有覆叠群}，它容许一个同时也是 Lie 群同态的光滑覆叠映射 $\pi\colon\widetilde{G}\to G$。
\end{theorem}

\begin{proof}
设 $\widetilde{G}$ 是 $G$ 的万有覆叠流形，$\pi\colon\widetilde{G}\to G$ 是相应的光滑覆叠映射。由\cref{exercise:ch04-38}，$\pi\times\pi\colon\widetilde{G}\times\widetilde{G}\to G\times G$ 也是光滑覆叠映射。

记 $m\colon G\times G\to G$ 和 $i\colon G\to G$ 分别为 $G$ 的乘法映射和取逆映射，并取 $\pi^{-1}(e)\subseteq\widetilde{G}$ 中任意元素 $\tilde{e}$。由于 $\widetilde{G}$ 是单连通的，覆叠映射的提升判据（\cref{proposition:appA-78}）保证映射 $m\circ(\pi\times\pi)\colon\widetilde{G}\times\widetilde{G}\to G$ 有唯一的连续提升 $\widetilde{m}\colon\widetilde{G}\times\widetilde{G}\to\widetilde{G}$，满足 $\widetilde{m}(\tilde{e},\tilde{e})=\tilde{e}$ 且 $\pi\circ\widetilde{m}=m\circ(\pi\times\pi)$：
\begin{align*}
\begin{tikzcd}[ampersand replacement=\&]
\widetilde{G}\times\widetilde{G} \ar{r}{\widetilde{m}} \ar{d}{\pi\times\pi} \& \widetilde{G} \ar{d}{\pi}\\
G\times G \ar{r}{m} \& G
\end{tikzcd}
\end{align*}
由于 $\pi$ 是局部微分同胚且 $\pi\circ\widetilde{m}=m\circ(\pi\times\pi)$ 是光滑的，由\cref{exercise:ch04-10}(a)可知 $\widetilde{m}$ 是光滑的。同理，$i\circ\pi\colon\widetilde{G}\to G$ 有光滑提升 $\widetilde{\iota}\colon\widetilde{G}\to\widetilde{G}$，满足 $\widetilde{\iota}(\tilde{e})=\tilde{e}$ 且 $\pi\circ\widetilde{\iota}=i\circ\pi$：
\begin{align*}
\begin{tikzcd}[ampersand replacement=\&]
\widetilde{G} \ar{r}{\widetilde{\iota}} \ar{d}{\pi} \& \widetilde{G} \ar{d}{\pi}\\
G \ar{r}{i} \& G
\end{tikzcd}
\end{align*}
我们在 $\widetilde{G}$ 中定义乘法与取逆为 $xy=\widetilde{m}(x,y)$ 与 $x^{-1}=\widetilde{\iota}(x)$，对所有 $x,y\in\widetilde{G}$。于是上述两图可改写为
\begin{align}
\pi(xy)=\pi(x)\pi(y),\label{align:ch07-03}\\
\pi(x^{-1})=\pi(x)^{-1}.\label{align:ch07-04}
\end{align}
接下来只需证明 $\widetilde{G}$ 在这些运算下构成群，因为那样 $\widetilde{m}$ 和 $\widetilde{\iota}$ 是光滑的，且\eqref{align:ch07-03}表明 $\pi$ 是同态，故 $\widetilde{G}$ 是 Lie 群。

首先证明 $\tilde{e}$ 是 $\widetilde{G}$ 中乘法的单位元。考虑映射 $f\colon\widetilde{G}\to\widetilde{G}$，定义为 $f(x)=\tilde{e}x$。则\eqref{align:ch07-03}蕴涵 $\pi\circ f(x)=\pi(\tilde{e})\pi(x)=e\pi(x)=\pi(x)$，故 $f$ 是 $\pi\colon\widetilde{G}\to G$ 的一个提升。恒等映射 $\operatorname{Id}_{\widetilde{G}}$ 是 $\pi$ 的另一个提升，且它与 $f$ 在一点处一致，因为 $f(\tilde{e})=\widetilde{m}(\tilde{e},\tilde{e})=\tilde{e}$，故覆叠映射的唯一提升性质（\cref{proposition:appA-77}(a)）蕴涵 $f=\operatorname{Id}_{\widetilde{G}}$，等价地，对所有 $x\in\widetilde{G}$ 有 $\tilde{e}x=x$。同理可证 $x\tilde{e}=x$。

接下来证明 $\widetilde{G}$ 中的乘法是结合的，考虑两个映射 $\alpha_L,\alpha_R\colon\widetilde{G}\times\widetilde{G}\times\widetilde{G}\to\widetilde{G}$，定义为
\begin{align*}
\alpha_L(x,y,z)=(xy)z,\qquad \alpha_R(x,y,z)=x(yz).
\end{align*}
反复应用\eqref{align:ch07-03}蕴涵
\begin{align*}
\pi\circ\alpha_L(x,y,z)=\bigl(\pi(x)\pi(y)\bigr)\pi(z)=\pi(x)\bigl(\pi(y)\pi(z)\bigr)=\pi\circ\alpha_R(x,y,z),
\end{align*}
故 $\alpha_L$ 和 $\alpha_R$ 都是同一映射 $\alpha(x,y,z)=\pi(x)\pi(y)\pi(z)$ 的提升。由于 $\alpha_L$ 和 $\alpha_R$ 在 $(\tilde{e},\tilde{e},\tilde{e})$ 处一致，它们相等。类似的论证表明 $x^{-1}x=xx^{-1}=\tilde{e}$，故 $\widetilde{G}$ 是群。
\end{proof}

\begin{exercise}\label{exercise:ch07-08}
通过证明 $x^{-1}x=xx^{-1}=\tilde{e}$ 来完成上述定理的证明。
\end{exercise}

我们还有如下唯一性结果。

\begin{theorem}\label{theorem:ch07-09}[万有覆叠群的唯一性]
对任意连通 Lie 群 $G$，万有覆叠群在如下意义下唯一：若 $\widetilde{G}$ 和 $\widetilde{G}'$ 都是单连通 Lie 群，且分别容许同时也是 Lie 群同态的光滑覆叠映射 $\pi\colon\widetilde{G}\to G$ 和 $\pi'\colon\widetilde{G}'\to G$，则存在 Lie 群同构 $\Phi\colon\widetilde{G}\to\widetilde{G}'$ 使得 $\pi'\circ\Phi=\pi$。
\end{theorem}

\begin{proof}
见\cref{exercise:ch07-05}。
\end{proof}

\begin{example}\label{example:ch07-10}[万有覆叠群]
\begin{enumerate}[(a)]
\item 对每个 $n$，映射 $\varepsilon^n\colon\mathbb{R}^n\to\mathbb{T}^n$ 由
\begin{align*}
\varepsilon^n(x^1,\dots,x^n)=\bigl(e^{2\pi\mathrm{i}x^1},\dots,e^{2\pi\mathrm{i}x^n}\bigr)
\end{align*}
给出，是 Lie 群同态和光滑覆叠映射（见\cref{example:ch07-04}(d)）。由于 $\mathbb{R}^n$ 是单连通的，这表明 $\mathbb{T}^n$ 的万有覆叠群是加法 Lie 群 $\mathbb{R}^n$。

\item \cref{example:ch07-04}(c)中描述的 Lie 群同态 $\exp\colon\mathbb{C}\to\mathbb{C}^*$ 也是光滑覆叠映射，故 $\mathbb{C}$ 是 $\mathbb{C}^*$ 的万有覆叠群。
\end{enumerate}
\end{example}

\end{document}